% =============================================================================
% Kapitel 2: Risiken, Failure Cases, Methodik und Meilensteine
% =============================================================================

\section{Risiken im industriellen 24/7-Betrieb}

Systeme, die ohne Unterbrechung laufen, sind einer Klasse von Problemen ausgesetzt,
die im Entwicklungsbetrieb selten sichtbar werden. Für ROSI lassen sich die
relevanten Risikokategorien wie folgt einteilen:

\begin{itemize}
    \item \textbf{Speichererschöpfung:} Puffer, Logs oder Datenbanken wachsen schleichend,
    bis Schreiboperationen systemweit fehlschlagen.

    \item \textbf{Stille Prozessausfälle:} Ein Teilprozess stürzt ab, das System läuft scheinbar
    weiter -- ohne Alarm und ohne Neustart.

    \item \textbf{Verbindungsabbrüche:} Datenbankverbindungen veralten nach Produktionspausen
    und führen beim nächsten Zugriff zu unkontrollierten Prozessabstürzen.

    \item \textbf{Kaskadeneffekte:} Die Verlangsamung einer Komponente blockiert über gemeinsame
    Queues und Shared Memory das gesamte System.

    \item \textbf{Blinde Flecken im Monitoring:} Vorhandene Überwachungskomponenten erfassen
    nicht alle kritischen Ressourcen, etwa flüchtige RAM-Dateisysteme oder interne
    Queue-Tiefen.
\end{itemize}

ROSI verfügt aktuell über kein externes Alerting. Probleme werden erst durch ausbleibende
Inspektionsergebnisse oder manuelle Beobachtung sichtbar. Die vorliegende Fallstudie
untersucht Fehler dieser Klasse, die im Dauerbetrieb entstehen und reproduzierbar
simuliert werden können.

\section{Themenblöcke und Failure Cases}

Die identifizierten Failure Cases werden in drei Themenblöcke gegliedert, deren
inhaltliche Ausrichtung nachfolgend auf übergeordneter Ebene dargestellt wird.
Die detaillierte Analyse der einzelnen Cases erfolgt in der Ausarbeitung.

\subsection{Block A -- Speicher- und Ressourcenverwaltung}

Forschungsfrage: Wie verhält sich das System unter begrenzten Ressourcen, wenn
Subsysteme überlaufen oder Prozesse unkontrolliert enden?

Untersucht werden Szenarien rund um den flüchtigen Bildspeicher als Single Point of
Failure sowie Timing-Konflikte im Shared-Memory-Management des Worker-Pools unter Last.

\subsection{Block B -- Systembeobachtbarkeit und Fehlertoleranz}

Forschungsfrage: In welchem Zeitraum erkennt das System eigene Prozessausfälle, und
welche Reaktion erfolgt?

Untersucht wird die Lücke zwischen der vorhandenen Prozess-API (State-Management,
Start/Stop) und der fehlenden Supervisor-Logik in der Laufzeitüberwachung. Ergänzend
werden konkrete Blind Spots des eingeführten Load Monitors analysiert.

\subsection{Block C -- Datenpersistenz, Retention und Kaskadeneffekte}

Forschungsfrage: Wie verhält sich das System bei unterbrochenen Datenbankverbindungen
oder unkontrolliert wachsenden Datenmengen, und welche Folgeausfälle entstehen?

Untersucht werden reproduzierbare Verbindungsabbrüche nach Produktionspausen sowie
Kaskadeneffekte durch blockierende Log-Queues, die zu einem systemweiten Stillstand
führen können.

\section{Methodik}

Die Untersuchung erfolgt am Referenzsystem im Labor Bad Hersfeld (Xubuntu 24.04 LTS,
AMD Ryzen 9 9950X, NVIDIA RTX 5090). Alle Failure Cases sind vollständig ohne echte
Kamerahardware simulierbar; ein vorhandener Simulator erlaubt reproduzierbare Szenarien
unter kontrollierten Bedingungen.

Pro Case werden konkrete Messgrößen erhoben und systematisch ausgewertet:

\begin{itemize}
    \item Time-to-Detection bei Prozessausfällen
    \item Inspektionsdurchsatz (Inspections/s) unter variierender Last
    \item Queue-Tiefe und Speicherfüllstand über die Zeit
\end{itemize}

\section{Meilensteine}

\begin{table}[h]
    \centering
    \renewcommand{\arraystretch}{1.2}
    \begin{tabularx}{\textwidth}{|l|l|X|}
        \hline
        \textbf{Meilenstein} & \textbf{Datum} & \textbf{Zwischenergebnis} \\ \hline
        M1 & 18. Juni 2026 & Theorierahmen erarbeitet, Testumgebung aufgesetzt, erste Cases gemessen \\ \hline
        M2 & 16. Juli 2026 & Alle Failure Cases gemessen, Rohdaten vollständig vorhanden \\ \hline
        M3 & 13. August 2026 & Auswertung aller Cases abgeschlossen, Ausarbeitung begonnen \\ \hline
        Abgabe & 27. August 2026 & Ausarbeitung (15 Seiten) und Präsentation (15 Minuten) fertig \\ \hline
    \end{tabularx}
\end{table}
